%=========================
% Plantilla LaTeX Artículo Académico
% Compilar con: pdflatex + biber + pdflatex + pdflatex
%=========================
\documentclass[11pt,a4paper]{article}

%-------------------------
% Idioma y tipografía
%-------------------------
\usepackage[english]{babel}
\usepackage{csquotes}
\usepackage[T1]{fontenc}
\usepackage{lmodern}

%-------------------------
% Márgenes y espaciado
%-------------------------
\usepackage[a4paper,margin=2.5cm]{geometry}
\usepackage{setspace}
\onehalfspacing

%-------------------------
% Matemáticas y símbolos
%-------------------------
\usepackage{amsmath,amssymb,amsthm}

%-------------------------
% Tablas y figuras
%-------------------------
\usepackage{graphicx}
\usepackage{booktabs}
\usepackage{caption}
\usepackage{subcaption}
\usepackage{changepage}

%-------------------------
% Enlaces y PDF
%-------------------------
\usepackage{hyperref}
\hypersetup{
  colorlinks=true,
  linkcolor=blue,
  citecolor=blue,
  urlcolor=blue,
  pdftitle={The Continuum Hypothesis: Sign-Game or Mathematics?},
  pdfauthor={Rodrigo Blanco Betes}
}

%-------------------------
% Listas, código y utilidades
%-------------------------
\usepackage{enumitem}
\usepackage{xcolor}
\usepackage{csquotes} % recomendado con biblatex

%-------------------------
% Bibliografía (biblatex + biber)
%-------------------------
\usepackage[
  backend=biber,
  style=apa,
  sortcites=true,
  url=true,
  doi=true
]{biblatex}
\DeclareLanguageMapping{english}{english-apa}
\addbibresource{refs.bib}

%-------------------------
% Metadatos del artículo
%-------------------------
\title{\textbf{The Continuum Hypothesis: Sign-Game or Mathematics?}}

\author{Rodrigo Blanco Betes}


\date{\today}

%-------------------------
% Entornos de teoremas (opcional)
%-------------------------
\newtheorem{theorem}{Theorem}
\newtheorem{lemma}{Lemma}

\begin{document}
\maketitle

%-------------------------
% Resumen y palabras clave
%-------------------------
\begin{abstract}


\end{abstract}

\noindent\textbf{Key words:} Continuum Hypothesis, Conceptual Structuralism, mathematical legitimacy, practice of mathematics.

%=========================
\section{Introduction}
\label{Intro}

Toward the end of the 19th century, Georg Cantor developed a theory of cardinality for infinite sets. Among his key results was the proof that the cardinality of the continuum strictly exceeds that of the natural numbers. This prompted him to ask whether any set could exist with a cardinality strictly between the two --- a conjecture that came to be known as the \textit{Continuum Hypothesis} (CH). The argument generalizes naturally: for any set $A$, one can show that the cardinality of its power set $\mathcal{P}(A)$ is strictly larger than that of $A$ itself. This observation gives rise to the \textit{Generalized Continuum Hypothesis} (GCH), which asserts that no set of intermediate cardinality can exist between $A$ and $\mathcal{P}(A)$. 

The work of \textcite{Godel1940} and \textcite{Cohen1966} showed it to be independent of the ZFC axioms of set theory, and therefore impossible to decide using the tools usually available to mathematicians in practice. This has led some authors to question whether the question admits a unique and meaningful answer at all. Among the most prominent is Solomon Feferman, who has argued (\citeyear{Feferman2011a}) that CH is an indefinite mathematical statement---the mathematical structures in which it could be evaluated being, in his view, not clear and intrinsically vague.

In this paper I would like to face this controversy from a different angle. In my view, one of the most crucial and delicate features of any meaningful mathematical statement is whether it amounts to more than a mere logical tautology. Almost any mathematical result can of course be read as a true statement within a given formal system---but the real significance of mathematics lies in whether its symbols connect to something beyond that system, whether they carry meaning that is not exhausted by the rules governing them. This is a central concern of the \textit{philosophy of mathematical practice}, a growing field that takes seriously the gap between formal proof and mathematical understanding. 

This paper approaches the question through what I call \textit{quasi-empirical structuralism}, a perspective that draws on two distinct but complementary commitments. The first is the quasi-empiricist philosophy of mathematics developed by \textcite{Lakatos1976}, on which mathematical knowledge grows through an iterative process of conjecture, refutation, and revision driven by the practice itself---not by the top-down imposition of formal systems. Underlying this view is a broader claim about mathematical meaning: that the significance of a mathematical statement cannot be read off from its position within a formal system alone, but depends on the uses to which it is put and the problems it is called upon to solve. This idea finds a striking expression in Wittgenstein:

\begin{adjustwidth}{1cm}{1cm}
\small
\begin{quote}
It is essential to mathematics that its signs are also employed in mufti (...). It is the use outside mathematics, and so the meaning of the signs, that makes the sign-game into mathematics. \parencite[V §2]{Wittgenstein1978}\footnote{Wittgenstein uses the phrase “employed in mufti” to mean that mathematical signs are used outside their formal mathematical setting—in ordinary contexts such as science, measurement, engineering, or everyday reasoning. In these contexts, the symbols appear without their \textit{formal mathematical uniform}, that is, outside the strictly formal system.}
\end{quote}
\end{adjustwidth}

%------------------------------------
The second commitment is a structuralist approach to mathematical ontology. On this view, mathematical statements are formulated within structures of objects whose relevant features lie entirely in the relations among them, not in any intrinsic properties the objects themselves might have. Most structuralists hold that accepting such a structure carries genuine ontological weight---a commitment that cannot be dissolved into pure logic.\footnote{See \textcite{Shapiro1997}.} This, however, raises an immediate difficulty. Given any consistent set of axioms, a model satisfying them can be shown to exist---which would seem to oblige the structuralist to count every logically coherent structure as a legitimate mathematical one. Most structuralists resist this conclusion, but the grounds for doing so are non-trivial: it is not enough to point to logical consistency alone. This brings us to a further and more fundamental question, one that will prove central to the argument of this paper---why do certain structures count as genuinely mathematical, while others remain merely formally possible?

This last point bears directly on the debate surrounding the Continuum Hypothesis. It is of course true that CH has a precise and well-defined meaning within the language of set theory. But well-definedness within a logical framework does not, in general, suffice for genuine mathematical meaningfulness---if it did, even the most elaborate and idle tautologies would qualify as mathematical claims.


If the structures in which CH is formulated are to count as genuinely mathematical, some further justification is required. This is precisely the question this paper seeks to address: whether the abstract structures involved in the formulation of CH have moved beyond what Wittgenstein would call a merely formal \textit{sign-game}---whether they connect to something that grounds their mathematical legitimacy. The framework developed here draws closely on Feferman's \textit{conceptual structuralism}, on which mathematical objects are positions within structures determined by our conceptual practices, and the legitimacy of a mathematical statement depends on whether those structures are sufficiently clear and grounded.

\section{The structuralists' approach: ontology, epistemology and categoricity}
\label{structuralims-ontology-epistemology}
 
An important philosophical approach that distances itself from formalism while engaging seriously with both epistemological and ontological questions is mathematical structuralism, developed most prominently by Shapiro and Resnik. As will become clear, structuralist approaches---including Feferman's conceptual structuralism---bear directly on debates about the determinacy of set-theoretic statements such as CH, particularly the question of whether there is a single, fixed set-theoretic universe.

One of the most paradigmatic developments within this philosophical movement is that of \textcite{Shapiro1997}. In this work, he defends a conception known as \textit{ante rem structuralism}, according to which mathematical structures exist independently of—and prior to—the systems of objects that instantiate them. On this view, what matters is not which particular objects occupy positions within a given structure, but rather the relations that define the structure and the ways in which structures relate to and extend one another.

Shapiro remains nevertheless open to alternative forms of structuralism, acknowledging them as broadly equivalent in practice even if less natural with respect to the actual development of mathematics. This includes Resnik's pattern-based (\citeyear{Resnik1981}) or non-ontological structuralism (\citeyear{Resnik2018NonOntological}). Shapiro himself is explicit on this point:

\begin{adjustwidth}{1cm}{1cm}
\small
\begin{quote}
I believe that the \textit{ante rem} option is the most perspicuous and least artificial of the three [eliminative, modal eliminative and \textit{ante rem} structuralism]. It comes closest to capturing how mathematical theories are conceived. Nevertheless, I do not mean to rule out the other options. Indeed, it follows from the thesis of structuralism that, in a sense, all three options are equivalent. As will be shown, each delivers the same “structure of structures.” \parencite[p. 90]{Shapiro1997}
\end{quote}
\end{adjustwidth}

The most important upshot of ante rem structuralism, for the purposes of this paper, is that structures are the fundamental primitives of mathematics, and that mathematicians characteristically reason in structural terms. Chapters 5 and 6 of \textcite{Shapiro1997} illustrate how deeply this perspective is embedded in the history of mathematics. This is evident, for instance, in \textcite{Dedekind1888}, as well as in \textcite{Bourbaki1950}, where “structure” is conceived as a collection endowed with one or more relations satisfying a specified set of axioms. The philosophical consequence is a broadly holistic constraint on mathematical meaning: a statement can only be considered determinate or meaningful relative to the structures in which it is defined. Applied to CH, this means that the question of its meaningfulness turns on whether the structures that model the languages in which it is formulated---including set theory---are themselves sufficiently grounded to confer that meaning.

Considerations around structuralism naturally raise the question of whether mathematical structures exist prior to objects, but also introduce a further “chicken or egg” problem: are structures prior to the axioms that characterize them, or are they instead determined through axiomatic specification? \textit{Ante rem} structuralism suggests that axioms merely describe independently existing structures, although the position is less clear on this issue than on the question of object-structure priority.

It is at this point where the epistemological problems regarding structures come into play. According to \textcite{Shapiro1997}, and much in line with \textcite{Resnik1981}, the fundamental way of coming to know a mathematical structure is through pattern recognition. The clearest case is the natural number structure, where the relevant pattern is the indefinitely iterable act of adding one element to a finite sequence:

\begin{adjustwidth}{1cm}{1cm}
\begin{quote}
  \small
Our subject learns that the sequence of sequences goes on indefinitely. For each sequence, there is a unique next-longest sequence, and so there is no longest sequence. The system of finite sequences is potentially infinite. Eventually, the subject can coherently discuss the structure of these finite patterns, perhaps formulating a version of the Peano axioms for this structure. We have now reached the structure of the natural numbers. \parencite[pp. 118-119]{Shapiro1997}
\end{quote}
\end{adjustwidth}

At this point we see the sharp line between pattern recognition and axiom systems. When more abstract mathematical structures, such as the set theoretical hierarchy, come into play, then the patterns leading to them may be more complex. At this point, the pattern involved leads to fundamentally an axiom system, such as $\mathrm{ZFC}$, and the structure is apprehended not directly but through what Shapiro calls implicit definition:

\begin{adjustwidth}{1cm}{1cm}
\begin{quote}
  \small
One way to understand a particular structure is through a direct description of it (...). Eventually, a properly prepared listener will understand that it is the structure itself, and not any particular instance of it, that is being described (...). We have here an instance of implicit definition, a technique familiar in mathematical logic. \parencite[p. 130]{Shapiro1997}
\end{quote}
\end{adjustwidth}

According to Shapiro, thus, relevant structures in mathematical discourse and axioms describing them need to stay at the same level, which in logical terms we understand as having structures determined by categorical axiom systems. However, crucial axiom systems for the practice of mathematics, such as $\mathrm{PA}$ for arithmetic and $\mathrm{ZF}$ for set theory, are not categorical in first order logic.\footnote{As categoricity of a theory implies completeness, any theory affected by Gödel's incompleteness theorem won't be categorical.} Accordingly, \textcite{shapiro1991foundations} argued that second order logic, where these theories are categorical (quasicategorical in the case of set theory), provides a suitable foundation for mathematics.\footnote{See \textcite[chapter 4]{shapiro1991foundations} for a proof of those categoricity results.} In chapter 8 of this book, Shapiro defends that, in pre-formal practice, mathematicians operate with concepts such as ``all subsets'', ``all natural numbers'', that fixes a unique interpretation for the axioms which enable them to distingish stardard from non standard models. Second-order logic, under full semantics, is precisely the formalism that is faithful to those informal semantic commitments, whereas first-order logic systematically fails to be.

This perspective, however, is subject to some important criticism. \textcite{Vaananen2015}---much in line with \textcite{Quine1941}--- argues that second-order logic depends on what he calls large domain assumptions, which come quite close to the meaning of the axioms of set theory. \textcite{Shapiro2012} even accepts this, but holds that it does not undermine his case for second order logic as a natural foundation of mathematics. They agree, for that reason, that second order logic and first order set theory can be threated as logically equivalent. However, from the perspective of the practice of mathematics, the latter is no doubt more prominent in the various foundationalist approaches to mathematics. 

Arguments that Shapiro gives related to mathematical practice are also problematic. For example he states (\citeyear{shapiro1991foundations}, p. 99) that important tool in mathematical practice, which is the concept of a minimal closure, cannot be formalized without appealing to second order logic: ``There is a straightforward argument, relying on the compactness theorem, that no collection of first-order formulas can successfully characterize any non-trivial minimal closure. This is evidence that first-order languages are inadequate.'' However, applying this argument to mathematical practice is not that straightforward, since in general practicioners tend to use minimal closures that are defined within a previously given structure. Take, for instance, the algebraic closure of $\mathbb{Q}$, which is one of the most paradigmatic examples. This set contains all the solutions to polynomical equations with coeficients in $\mathbb{Q}$, which can be defined with parameters over $\mathbb{Q}$ in the structure $\mathbb{C}$, using the language of first-order theory of rings plus a really weak version of first-order set theory. More generally, in a given structure, most closures are definable by formulas with parameters or describable as the intersection of definable sets. Shapiro's compactness argument establishes a genuine limitation of first-order logic in the abstract, but it does not map cleanly onto the way closures are actually used.

%-------------Distinction: open-ended vs determined structures----------------
In this discussion---the priority of axioms and patterns, or that of structures over them---not all the structuralists agree Shapiro. \textcite[Chapter 7]{Resnik1997} defends that a mathematical theory, given by axioms, merely describes a structural pattern, without necessarily determining a unique and privileged structure that models that theory.

A clear and paradigmatic case of this approach in mathematical practice is category theory. With a remarkably concise set of axioms---objects, morphisms, composition, and identity---the theory applies across a wide variety of mathematical domains, including groups, topological spaces, sets, relations, partially ordered sets, and differentiable manifolds (see section \ref{cat-theory}). What matters in practice, however, are the particular collections that instantiate these axioms, rather than whether there exists a fundamental and privileged structure that realises them. This is precisely the spirit in which \textcite{maclane1998categories} develops the theory: category theory is presented as a unifying language for mathematical practice, one whose power lies in its ability to reveal structural analogies across disparate areas, not in grounding mathematics in a single canonical structure. MacLane is notably cautious about foundational overreach---he treats the axioms of category theory as a flexible and powerful descriptive tool rather than as the implicit definition of a unique intended model.\footnote{Similarly, \textcite{Feferman2013Category} argues that the strength of category theory lies precisely in what MacLane emphasised---its role as a structural lingua franca across mathematical domains---rather than in any foundational pretension. He is also critical of the notion of a ``category of all categories'': beyond the coherence problems such a universal structure faces, he argues that explaining what a category is already presupposes a prior notion of collection and operation, so category theory cannot serve as a logically autonomous foundation.}

The category theory case thus illustrates, in a particularly vivid way, the broader tension between the two modes of structural commitment identified above. Neither of these two approaches---the priority of structures, or the priority of axioms and patterns---is unanimously adopted in mathematical practice, and I do not think one should be universally privileged over the other. Rather, the two capture genuinely different modes of mathematical commitment, and the distinction between them is philosophically consequential. 

When mathematicians grasp a structure prior to the axiom system that characterises it---formulating axioms only later as fundamental statements about something already implicitly understood---I say they are \textit{committing to a mathematical structure}. This aligns with the spirit of \textit{ante rem} structuralism. In the opposite case, when mathematicians proceed in the reverse direction---beginning from an axiom system before having a clear structure in view, as has been increasingly common since the formalist turn---I say they are \textit{committing to an axiom system}. This is to prioritise the pattern over any particular structure that realises it. Crucially, when the axiom system in question is non-categorical, such a commitment does not determine a unique structure: many non-isomorphic models may satisfy the same axioms, and no principled grounds exist for privileging one over the others. By Gödel's completeness theorem, any consistent axiom system has a model; assuming mathematicians commit only to systems they take to be consistent, axiomatic commitment may therefore be identified with commitment to the entire class of structures satisfying those axioms. I refer to this class, for the remainder of the paper, as an \textit{open-ended structure}.

The distinction bears directly on the question of CH. Committing to a mathematical structure implies accepting a determinate mathematical reality in which statements formulated in the corresponding language have definite truth values. Committing to an axiom system, by contrast, implies only acceptance of a pattern---one that may be realised across a range of structures, leaving statements such as CH without a determinate answer.

%=========================
\section{Stepping away from formalism: mathematical practice and metodology}
\label{history-section}

As we have seen, structuralism provides important ontological and epistemological insights into mathematics, particularly through the notions of commitment to mathematical structures and axiom systems. Our concern here, however, is not merely the nature of such commitments, but the reasons mathematicians adopt particular structures or axioms in actual mathematical practice.

Philosophical approaches to mathematics have evolved dramatically over the course of the twentieth century. At the beggining of this century, both the mathematical world—strongly influenced by Hilbert’s formalism—and several philosophical approaches to mathematics, such as logical positivism, located the essence of mathematics in formal and logical methods. From the perspective of formalism, mathematical statements were understood as deductions from a given set of axioms, and any consistent theory formulated in a logical language was regarded as sufficient to define a corresponding piece of mathematical reality—see, for instance, \textcite{hilbert1971geometry}. Within logical positivism, as represented for example by \textcite{carnap1937syntax}, mathematical sentences were regarded as analytic, in the sense that they say nothing about the empirical world and function instead as tools for clarifying the structure of scientific language.

In the latter half of the twentieth century, new philosophical conceptions of mathematics emerged that emphasized practical applicability, heuristic reasoning, and the essentially informal character of everyday mathematical work. As pointed out in the \nameref{Intro}, this shift marked the advent of the so-called \textit{philosophy of mathematical practice}. One of its earliest and most influential proponents was Imre Lakatos, who, in his work \textit{Proofs and Refutations}, strongly criticized the formalist approach of earlier schools:

\begin{adjustwidth}{1cm}{1cm}
\small
\begin{quote}
Since, [according to logical positivism], informal mathematics is neither `tautological' nor empirical, it must be meaningless, sheer nonsense. The dogmas of logical positivism have been detrimental to the \textit{history and philosophy of mathematics}. \parencite[p. 2]{Lakatos1976}
\end{quote}
\end{adjustwidth}

Lakatos’ alternative approach is known as quasi-empiricism, as he argued that mathematics develops in a manner analogous to the empirical sciences: one begins with a conjecture and subsequently seeks a proof. Although such a proof may initially appear convincing, it can generally be decomposed into lemmas and implicit assumptions, which may themselves be vulnerable to refutation. When such refutations arise, the original conjecture is not simply discarded but rather refined—whether by restricting its scope, revising underlying definitions, or modifying auxiliary assumptions.\footnote{This idea is extensively developed through the analysis of Euler's formula, which states that for convex polyhedra with $V$ vertices, $E$ edges, and $F$ faces, $V - E + F = 2$. See \textcite[pp. 10–41]{Lakatos1976}.}

According to Lakatos, this perspective remains meaningful even in the context of the increasing rigor of nineteenth-century mathematics, particularly with the formalization of analysis:

\begin{adjustwidth}{1cm}{1cm}
\small
\begin{quote}
In the early nineteenth century the flood of counterexamples brought confusion. Since proofs were crystal-clear, refutations had to be miraculous freaks, to be completely segregated from the indubitable proofs. Cauchy’s revolution of rigour rested on the heuristic innovation that the mathematician should not stop at the proof: he should go on and find out what he has proved by enumerating the exceptions, or rather by stating a safe domain where the proof is valid.\parencite[p. 55]{Lakatos1976}
\end{quote}
\end{adjustwidth}

From the perspective developed here, Lakatos draws a clear distinction between mathematical reasoning and purely formal calculation. Rather than grounding mathematics in formal derivations alone, he emphasizes its problem-solving power: definitions, entities, and assumptions are justified insofar as they contribute to the advancement of inquiry. Another central feature of his view is the iterative character of mathematical reasoning. The emergence of a refutation or the identification of a flaw does not terminate the process; instead, it prompts the introduction of revised assumptions, refined definitions, or new intermediate results, which may in turn generate further tensions and revisions. This iterative character is central to Lakatos’ account, according to which mathematics advances not through the accumulation of final proofs, but through a dynamic process of conjecture and criticism. As he puts it (\citeyear{Lakatos1976}, p.5), “mathematics [\dots] does not grow through a monotonous increase of the number of indubitably established theorems but through the incessant improvement of guesses by speculation and criticism, by the logic of proofs and refutations”.

%--------Inserting MacLane-------------------
Lakatos' approach illuminates not just how mathematics is justified, but how it develops---through problems, failures, and the iterative refinement of concepts. This raises a question central to the present paper: how do mathematical structures and axiom systems actually come into being, and why are they introduced when they are? This is a question that, in a similar vein to Lakatos, much oriented to practice, \textcite[36]{maclane1986mathematics} further explores:

\begin{adjustwidth}{1cm}{1cm}
\small
\begin{quote}
The genesis of the more complex mathematical structures tends to take place within Mathematics itself. [There] are a variety of processes which may generate new ideas and new notions. We list a few of these processes in tentative form for further refinement after our more detailed studies. \parencite[p. 36]{maclane1986mathematics}
\end{quote}
\end{adjustwidth}

MacLane's list includes \textit{conundrums} (the pressure exerted by hard unsolved problems), completion, invariance, common structure (the recognition of structural analogies across different areas), generalisation, abstraction, axiomatisation, and analysis of proof. What is striking about this catalogue is how closely it mirrors the Lakatosian picture: mathematical structures do not arrive fully formed, imposed from above by logical fiat, but are drawn out by the demands of practice. MacLane makes this developmental dimension explicit later in the same work, arguing that mathematics emerges from human experience and scientific practice and grows through successive cycles of formalisation and generalisation:

\begin{adjustwidth}{1cm}{1cm}
\small
\begin{quote}
Mathematics, in our description, rests on ideas and problems arising from human experience and scientific phenomena and consists in many successive and interconnected steps in formalizing and generalizing these inputs. \parencite[449]{maclane1986mathematics}
\end{quote}
\end{adjustwidth}

%--------------Externally applicable structures
The reference to human experience and science is crucial here. MacLane's account of how mathematics develops has the shape of a mathematical induction. There is a base case: a set of notions directly grounded in practical and scientific engagement with the physical world---measurement, counting, spatial reasoning, the mathematics demanded by physical problems. From this foundation, the iterative processes MacLane catalogues---generalisation, abstraction, the recognition of common structure, the analysis of proof---allow the totality of mathematical knowledge to grow, each step building on what came before while progressively loosening its ties to the physical ground. What results is a body of mathematics that is largely self-sustaining and internally driven, yet never fully detached from the base that gave rise to it.

This duplicity---between mathematics as an autonomous discipline pursued for its own sake and mathematics as something that cannot afford to lose contact with the external world---is captured with remarkable precision by \textcite[pp. 332--337]{Poincare1897rapports}:

\begin{adjustwidth}{1cm}{1cm}
\small
\begin{quote}
And above all, those devoted to mathematics find in it pleasures analogous to those given by painting and music (\dots). That is why I do not hesitate to say that mathematics deserves to be cultivated for its own sake, and that theories which cannot be applied to physics should be esteemed just as highly as the others (...). But the pure mathematician who would forget the existence of the external world would be like a painter who knew how to harmoniously combine colors and forms, but lacked models. His creative power would soon run dry. 
\end{quote}
\end{adjustwidth}



How could we apply this insights to the perspective developed here, that assumes a structuralist approach to ontology in mathematics? 

\textbf{1: Which are the \textit{original} mathematical structures}

This is fundamentally the idea of a ``base case'' in the collection $\mathcal{S}$, which in my view should contain structures whose legitimacy is grounded directly in their indispensability---or at least their deep entrenchment---in scientific and practical engagement with the physical world. This does not require accepting the indispensability argument of \textcite{quine1960} and \textcite{putnam1971} in its full ontological strength. The weaker claim defended here is simply that the most successful structures in modelling the physical world and ubiquitous in scientific practice should belong to any reasonable conception of mathematics. A structure satisfying these conditions relative to extra-mathematical practice will be called an \textit{externally applicable structure}.

I propose to understand it through the notion of $\mathcal{S}$, the collection of (possibly open-ended) mathematical structures to which mathematical practice is committed. 

One I defend to be of the most important insights of Lakatos’ quasi-empiricism, 

\textbf{2: how do mathematical structures appear from others}

is the fact that $\mathcal{S}$ is closed in the following sense: a mathematical structure $X$ should be also included in that collection in case it satisfies at least one of the following conditions. Such a list will probably allways be tentative:

\begin{enumerate}

    \item \label{condition1} \textbf{Problem-solving power:} $X$ provides solutions to previously existing, well-defined, and relevant problems formulated within structures in $S$.

    \item \label{condition2} \textbf{Explanatory unification and simplification:} $X$ reveals that phenomena arising in structures in $S$, previously regarded as disparate, are in fact manifestations of a single underlying principle.

    \item \label{condition3} \textbf{Generative capacity:} $X$ gives rise to new and fruitful lines of inquiry within $S$.
\end{enumerate}

To avoid redundancy, $X$ should moreover be minimal or sufficiently simple relative to the problems solved, phenomena unified, or research lines generated. I return to this criterion in Section \ref{simplicity}. These three criteria capture central ways in which mathematical structures prove fruitful in practice, a point illustrated through the historical examples discussed in Section \ref{quasi-emp-structuralism}.

Note that this, as in Lakatos’ approach, proposes an inductive, iterative pattern, in which new structures and definitions emerge from previously accepted ones.  This is a point on which greater subtlety is required. 



The perspective developed here may therefore be understood as a form of \textit{quasi-empirical structuralism}: a structuralist conception of mathematics in which the legitimacy and development of mathematical structures are closely connected to their explanatory, problem-solving, and generative role within both mathematical and scientific practice. I think this conception is also faithful to the \textcite{Wittgenstein1978} passage mentioned in the \nameref{Intro}, insofar as it requires mathematical structures to be applicable to problems, phenomena, and areas of inquiry beyond their own immediate definition, and ultimately beyond mathematics itself.

What quasi-empirical structuralism defends, then, is a kind of methodological well-foundedness for mathematical structures. If $S_1$ and $S_2$ are mathematical structures,\footnote{I insist that such structures may be open-ended, meaning—see the end of section \ref{structuralims-ontology-epistemology}—that they may consist simply of collections of axioms accepted in practice.} I write 

$$S_1 \leadsto S_2 $$

to indicate that $S_2 \in \mathcal{S}$ because it fulfulls conditions \ref{condition1}–\ref{condition3} with respect to $S_1 \in \mathcal{S}$. In this case I would say, in a way that is natural regarding those conditions, that $S_2$ is \textit{internally fruitfull} with respect to $S_1$. Then for any structure $S \in \mathcal{S}$, there is a sequence 

$$ S_0 \leadsto ... \leadsto S_n = S $$ 

where $S_0$ is an externally applicable structure. What concerns us here, then, is whether the structures appearing in such a sequence, in the particular case where $S$ is involved in the formulation of $CH$, are open-ended or definite.\footnote{In section \ref{simplicity} I defend that, if $S \leadsto S'$ and $S$ is open-ended, then $S'$ must also be open-ended.} More precisely, the question is whether the fruitfulness of those structures in mathematical practice derives primarily from commitment to definite mathematical objects, or instead from the continued development of inferential and axiomatic patterns that remain open-ended while guiding mathematical activity.


\section{A practice-based, quasi-empirical structuralism}
\label{quasi-emp-structuralism}

In this section I examine several historical examples drawn from mathematical practice. The first objective is to argue that the iterative component of quasi-empirical structuralism has been present at various moments in the historical development of mathematics. The second is to defend the claim that structures such as $\mathbb{N}$, $\mathbb{R}$, and $\mathbb{C}$ have proved fruitful in mathematical practice precisely through commitment to their structural form, making it difficult to regard them as merely open-ended structures.

%---------------Complex numbers--------------------------
The early stages of this development can be traced back to the sixteenth century. Although the idea of negative square roots had already been mentioned before, it was Rafael Bombelli, in his 1572 work L’Algebra, the first to systematically use of what we now call complex numbers.\footnote{Most of the events mentioned here are important facts in history of mathematics, where I am following texts on the topic such as \textcite{Nahin2010} and \textcite{Christen2018}.} Bombelli was working on the solution of the reduced cubic equation, an equation of the form $x^3 + px + q = 0$. Earlier, Cardano had found the following formula to solve such equations:

\begin{equation}
x = \sqrt[3]{-\frac{q}{2} + \sqrt{\left(\frac{q}{2}\right)^2 + \left(\frac{p}{3}\right)^3}} 
+ \sqrt[3]{-\frac{q}{2} - \sqrt{\left(\frac{q}{2}\right)^2 + \left(\frac{p}{3}\right)^3}}
\end{equation} 

Cardano’s difficulty lay in interpreting this formula when negative square roots were involved. In fact, he concluded that the formula was not solvable “in general.” Bombelli decided to step a little bit further, introducing the symbol $pdm$ to represent $\sqrt{-1}$, thereby providing a clear precursor to the imaginary unit.\footnote{This stands literally for \textit{piú di meno}, which mean “plus or minus” in Italian—see \textcite{Christen2018} for more details. He arrived at many of the operational properties of complex numbers, such as $pdm \cdot pdm = -1$.} In this way, he was able to obtain new real roots of certain equations of the given form, which Cardano had been unable to.

Progress in this area accelerated in the following years. While Descartes showed little interest in such numbers, referring to them pejoratively as “imaginary,” Leibniz used them in a manner similar to Bombelli, as intermediate steps in solving equations. Euler was probably the first to use them systematically: he introduced the notation $i$ for the imaginary unit and employed it in many formulas. Perhaps the most well known is Euler's formula, which states that for any real phase $\alpha$, $e^{i\alpha} = \cos \alpha + i \sin \alpha$. Another notable application is the factorization of fourth-degree polynomials with no real roots (such as $x^4 + 1$). Later, Gauss introduced the complex plane, standardized the notation, and firmly established the acceptance of complex numbers, using them in many important results, most notably in the fundamental theorem of algebra. Since then, complex numbers have become one of the cornerstones of mathematics, and the set of complex numbers is now regarded as a mathematical structure in its own right.

What were the reasons for such a shift? In my view, complex numbers constituted the simplest mathematical structure extending the real line while exhibiting three fundamental features, which had traditionally—albeit in different forms throughout history—been regarded as criteria of legitimacy.

First, complex numbers demonstrated significant problem-solving power, as seen in the factorization and solution of higher-degree polynomials. Second, they provided explanatory unification: results such as the fundamental theorem of algebra, or the unification of trigonometric and exponential functions through Euler's formula, illustrate this point, and many further examples could be given. Finally, they enabled mathematicians to simplify and extend existing formulations to such an extent that new areas of application emerged, through connections with geometry, analysis, algebra, number theory, and beyond.

These features of structures such as the complex numbers are, in my view, what distinguish mathematical reasoning from mere symbol manipulation. 

For example, the fact that $(\mathbb{C}, +, \cdot)$ is an algebraically closed field indicates that it is the simplest structure satisfying these criteria relative to $\mathbb{R}$. Moreover, complex numbers play a crucial role in simplifying proofs across many areas of mathematics. For these reasons, if a given conception of the real line belongs to $S_n$, then the corresponding structure of complex numbers naturally belongs to $S_{n+1}$. For instance, a purely geometric conception of the real line—belonging to $S_0$—would correspond to a geometric conception of the complex numbers in $S_1$. By contrast, the more standard set-theoretic conception of the real line would correspond to a structure of complex numbers belonging to $S_2$, as argued in the previous section.


\textbf{Commitment to structures:} Our intuition clearly leans toward committing to $\mathbb{C}$ as a structure. This is really different from \textcite{Godel1947} ``intuition'' of set theoretical structures. 

%=========================
\section{The natural numbers: a structure well-determined by practice?}
Externally applicable structures: natural numbers, rationals, reals, geometry, ...

Determinateness of this structure: non standard models of $\mathbb{N}$ in practice. Real-life applicability: time tending to infinity, divisibility of matter, computer programs, \dots

Historical point: in contrast to set theory, axioms were given much later than the usage of these structures in practice.

\textcite{maclane1986mathematics} defends that the structure of natural numbers is based in practical neccesities.

\begin{adjustwidth}{1cm}{1cm}
\small
\begin{quote}
Other approaches restrict attention to those real numbers which are explicitly “constructible” in some specific sense, perhaps by means of recursive functions of natural numbers (...). In my view, such requirements of constructivity tend to emphasize arithmetic at the expense of geometry. Since the real numbers should subsume both arithmetic and geometric aspects of mensuration, these practical uses seem better formalized by the standard axioms for the field of real numbers. \textcite[107]{maclane1986mathematics}
\end{quote}
\end{adjustwidth}

Categoricity arguments around $\mathrm{PA}$. \textcite{Ferreiros2023ConStructuralism} says it is ``easier'' to proof the categoricity of $\mathrm{PA}$. $\mathrm{PA}$ standard model is formalizable within $\mathrm{ZFC}$? Arguments of \textcite{shapiro1991foundations} around second order logic are much better here: it is much easier to argue that second order logic better captures the informal language of mathematics in the case of arithmetic than in the one of set theory.
 

\section{Simplicity and naturalness in mathematical practice}
\label{simplicity}

\textbf{1. Goodman (and early Quine): attempts at formal measures of simplicity}

The closest thing to a genuinely systematic, quasi-mathematical account is due to Nelson Goodman (often in dialogue with W. V. O. Quine).

Goodman tried to define simplicity as “systematization” of a theory—roughly, how economical its conceptual scheme is.
He proposed axiomatic measures based on the logical structure of predicates (arity, definability, etc.).
The idea: a theory is simpler if it uses a less expressive (but still adequate) conceptual basis.

This is one of the very few attempts that looks like a formal calculus of simplicity.

Limitations:

Highly language-dependent (the “language variance” problem).
Hard to apply to real mathematics or science.
Captures only one dimension of simplicity (logical/conceptual).

Still, if you want something systematic in the strict sense, Goodman is essential.

\textbf{2. Popper, Peirce, and “simplicity = testability”}

Another semi-systematic tradition:

Karl Popper: simpler theories have greater empirical content / falsifiability.
C. S. Peirce: simplicity linked to ease of testing and elimination.

This gives a methodological account, not a structural one:

Simplicity = how strongly a theory constrains possible observations.

Problem:
It often misclassifies cases (Goodman’s famous counterexamples), so most philosophers think this captures one aspect rather than the whole story.

\textbf{3. Information-theoretic approaches (Kolmogorov complexity, MDL)}

More recent and mathematically precise:

Simplicity = shortest description length.
Formalized via Kolmogorov complexity or Minimum Description Length (MDL).

These are probably the closest modern analogues of a general theory:

A theory is simple if it compresses the data well.
Widely used in statistics, machine learning, and some philosophy of science.

Caveat:
They depend on a choice of language/model class (though only up to a constant), and they don’t fully capture “naturalness” in the mathematician’s sense.

Restricting from $\Pi_1^1-\mathrm{CA-}0$ to $\mathrm{ATR}_0$ does not significantly reduce Kolmogorov complexity at the level of axioms, but it does significantly reduce proof-theoretic strength and expressive/compressive power.

\textbf{5. Lakatos and mathematical practice (quasi-empiricism)}

Imre Lakatos doesn’t define simplicity formally, but gives something arguably deeper:

Mathematics evolves through proofs, counterexamples, and concept refinement.
“Natural” definitions are those that survive this dialectical process.

So instead of a formula:

Naturalness = robustness under mathematical practice.

This is influential if you’re thinking about how mathematicians actually judge elegance or naturalness.

\textbf{6. Contemporary philosophy: pluralism about simplicity}

Modern consensus (roughly):

There is no single notion of simplicity.
Instead, there are multiple dimensions:
ontological (few entities),
syntactic (short descriptions),
conceptual (few primitives),
explanatory (unification),
statistical (compression / likelihood).

Some philosophers even treat “simplicity” as a placeholder for a bundle of theoretical virtues rather than something reducible to one principle.


\section{How CH and set theory fit in quasi-empirical structuralism}
\label{section-CH-case}

One such kind of problem is known as the Weil Conjectures. Named after the mathematician that introduced them in 1949, consist of four deep hypotheses on algebraic varieties defined over finite fields. These conjectures were crucial for the development of modern algebraic geometry.  

The Weil conjectures can be formulated in an extremely weak axiom system, such as second order arithmetic, or even PA with minor additions. Consequently, they cannot be considered illegitimate for being far ahead of down-to-earth mathematics. However, proving them requires much more mathematical power. The conjectures were ultimately proven by Pierre Deligne in (1974), using tools previously developed by Grothendieck, using the concept of a Grothendieck Universe, which requires ZFC plus the existence of a strongly inaccessible cardinal. \textcite{McLarty2020} demonstrated that developing such tools requires a far weaker axiom system, one that can be formulated within finite-order arithmetic. In any case, these systems, which McCarty refers to as $MC$, require for each natural number $n$, the iterated power set $P^n(\mathbb{N})$ to exist.  $MC$ systems are the weakest systems where the proof of Weil’s conjecture using the tools Grothendieck provided can be carried out.  

The key point here is that within any such $MC$ system, CH is a perfectly definable statement. Consequently, if the proof of Weil’s Conjectures is accepted, we should therefore commit (for its internal fruitfulness regarding these conjectures) to the legitimacy of such axiomatic systems, and the corresponding structures that model them, and therefore to the legitimacy of CH.10 Up to this point, if one wanted still to defend the illegitimacy of CH, there would remain fundamentally the following options: 

\begin{enumerate}
  \item Not accepting Deligne’s proof, considering it is not based on full-predicative reasoning.  

  \item Returning to a more radical version of conceptual structuralism, considering only legitimate structures the ones that can be directly externally applied to science. 
\end{enumerate}

One such question is Fermat's Last Theorem. \textcite{Feferman1999} defended the idea that questions like this one would most likely be provable using tools in weaker systems, such as PA. In this particular case, however, the proof required much stronger tools so that, even in the most conservative reconstructions, can only be formulated within sixth-order arithmetic—see \textcite{McLarty2020}.

\textbf{Important remark:} In the case of sixth order arithmetic, it is not true either that the structures involved are prior to the axiom systems involved. In this regard, we can still stand against the commitment to the structures involved.


\section{The approach to set theory}
\label{set-theory}

Recall that the standard model of arithmetic can be formalized using the ZF axiom system. Therefore, the consequences of these axioms could not contradict a true result in arithmetic unless the axioms were found to be inconsistent, which seems highly unlikely. For this reason, even a mathematician only committing to the definiteness of arithmetic would find these axioms interesting. They could use the ZF axioms as a useful metamathematical tool and have reasons to believe in the consequences of the axioms when applied to arithmetic. But this would not lead them to commit to the full set-theoretic hierarchy, nor to defend that the concept of an arbitrary set is definite.\footnote{Even acknowledging that ZF is a consistent theory does not imply the aforementioned commitments. A Skolem-type numerable model could suffice for proving the consistency of ZF set theory intuitively.}

\textbf{Statements in arithmetic such as Goodstein's theorem were proven through axiom systems whose scope is beyond arithmetic, in particular transfinite ordinal induction. So it is not that unnatural to understand higher-order assumptions to refer in practice to arithmetic.}

If $L$ is a large cardinal axiom, the theories $ZFC+L$ and $PA + \bigcup_{n < \omega} Con((ZFC + L)\upharpoonright n)$ are mutually interpretable, being the latter one framed in the language of arithmetic. This means that the logical consequences of $ZFC + L$ in the language of set theory can be understood as arithmetical sentences through a well chosen bijection. That fits well with the formalist approach to the metamathematical results, where set theoretical sentences and deductions are understood as the object of study, while the core of the reasoning is done in the metatheory of formal finitary logic. This in fact means that, even if accepting the fruitfulness of set theoretical considerations, specially regarding consistency proofs, that reasoning can be reduced to the language of arithmetic. So, in the end, the only mathematical structure one has to commit to in this case is the natural number structure, without need of a freestanding and platonic universe of sets.

Is there a mathematical commit to $ZFC + (V = L)$ or to $ZFC + MM$, which stands for Martin's Maximum Axiom, presented in \textcite{ForemanMagidorShelah1988}?

The whole cumulative hierarchy (7) is intended (or imaged) to produce all the sets that we need for Mathematics. Though it does not produce any ``universal set'' of all sets, it does produce sets of all finite sizes as well as many infinite sets. \textcite[362]{maclane1986mathematics}

\begin{adjustwidth}{1cm}{1cm}
  \small
  \begin{quote}
Within ZFC set theory, the cumulative hierarchy can be formally constructed. Before axioms are at hand, the hierarchy is a Platonic myth, clearly visible only to those with a sixth sense for sets. \textcite[373]{maclane1986mathematics}
  \end{quote}
\end{adjustwidth}

\begin{adjustwidth}{1cm}{1cm}
  \small
  \begin{quote}
Proof in Mathematics is both a means to understand why some result holds and a way to achieve precision. As to precision, we have now stated an absolute standard of rigor: A Mathematical proof is rigorous when it is (or could be) written out in the first order predicate language $\mathcal{L}(\in)$ as a sequence of inferences from the axioms ZFC, each inference made according to one of the stated rules.\textcite[373]{maclane1986mathematics}
  \end{quote}
\end{adjustwidth}

We have the sequence: 

$$ (\mathbb{N}, 0, S), (\mathbb{R},\tau_u, \beta), (\mathbb{C}, ...) \leadsto \mathrm{ZFC} $$

It is not a matter of evaluating how easy it is to identify standard models in set theory \parencite[p. 16]{Koellner2022}, but of acknowledging that standard models are not used in practice. In general, all needed is a countable transitive model, and that can be carried out strictly in the metatheory---see \textcite[p. 8]{Kunen2011}. 

Practice in set theory commits to structures or to axioms? It is interesting to note that, differently from other cases in mathematics, the development of axioms in set theory was prior to stablishing the set theoretical hierarchy---what Von Neumann did in (1920) to explain in which sense the axioms did make sense, to avoid paradoxes and to give an structured version of the universe of sets. In other words, in this case the structure was introduced \textit{ad hoc} for the axioms, while in other cases (such as arithmetic) it was the other way around.

In relative consistency proofs we do not need to beleive that the axioms of set theory are consistent: we take it as an assumption of the result.

In this regard, the distinction between mathematics and metamathematics is crucial. Axiomatic set theory often functions in this metamathematical role; it provides results about the proving power of specific axiom systems for formalized mathematical languages. A paradigmatic example of a metamathematical reasoning about set-theoretical concepts is found in independence results. These results show the impossibility of logically proving certain statements—most notably, CH—from the axioms of a theory like ZFC. Crucially, these independence proofs do not rely on the abstract set-theoretical structures whose status they investigate. Within the proofs, these structures appear merely as formulas in a formal language. The actual reasoning occurs from an external framework, known as the metatheory. In general, the metatheory is taken to be a much weaker system than ZFC, often requiring only finitistic tools. For example, we see in \textcite{Kunen2011}, p. 11: “All we assert as being `really true' is elementary finitistic reasoning: this is called the metatheory. Using this reasoning, we can explain what a logical formula is and what a formal proof is (…). Then, all discussions about infinite and uncountable sets take place within ZFC, which is called the formal theory”. 

Therefore, even if one accepts —as I do— that this metamathematical reasoning constitutes a crucial and fruitful part of current mathematical practice, its success legitimates the weaker tools it employs (axioms or concepts from logic, finitistic reasoning, forcing, etc.), not the abstract set-theoretical structures themselves. 

Are abstract set theoretical structures internally fruitful with respect to structures in other areas of mathematics? In this regard, probably the most important success of set theory has been its full integration as a unifying language and a conceptual framework for all of mathematics. The concept of set provides a single but powerful tool to define nearly all relevant mathematical objects, functions, numbers, spaces, and more. This, of course, demonstrates a significant explanatory fruitfulness that can be attributed to set theory, offering a deeply integrated language across mathematical areas. However, this success is linguistic and conceptual, not foundational. It does not require committing to the higher levels of the cumulative set hierarchy, neither to the full strength of all axioms in set theory. It is true, for example, that the Axiom of Choice is often used in practice, but it is not clear whether this is usually done with its full strength or referring to weaker versions of it (such as Dependent Choice).

With respect to the power set axiom, consider, for instance, the field of topology. Topology focuses on the study of topological spaces—structures defined as a set together with a collection of its subsets satisfying certain conditions. The power set axiom seems be inherent in the very definition, as the topology on a set is defined as subset of the power set, . Nevertheless, I would argue this is interpreted in a schematic way by the working topologist. In practice, one reasons about any given family of subsets of that fulfills the specific rules. Theorems are proven about this structure itself. This schematic reasoning about a family of subsets if fundamentally different from needing to quantify over a pre-existing, completed totality . A potential objection could be the definition of the discrete topology, where the space is formed through the entire power set. However, in most cases practical applications, this topology is used on countable sets. When it is applied to uncountable sets, it is typically within the realm of abstract set-theoretical considerations themselves—precisely the context whose necessity for more down-to-earth applications we are questioning. One could concede that a topologist, for simplicity, assumes they can always form the power set of any set they encounter in practice. However, for this purpose a far weaker system that ZFC, like the previously given MC, could suffice.

\textbf{Set theory is commitment to axioms: the practicioner finds it useful to assume that some properties of the sets he is using are going to hold. For example, that for any set $A$, there is a set $P(A)$ with the property of containing all the sets $B$ such that $B \subseteq A$. But that does not mean that they commmit to an univocally determined universe of sets.}

%=========================
\section{Category theory}
\label{cat-theory}

The axiomatic basis of this theory is the definition of a category $\mathcal{A}$, which consists of the following elements:

\begin{enumerate}
  \item A collection $Obj(\mathcal{A})$, whose members are called the objects of $\mathcal{A}$.
  \item For every pair of objects $A, B \in Obj(\mathcal{A})$, a collection of arrows (or morphisms) from $A$ to $B$. If $f$ is such an arrow, we write $f: A \rightarrow B$.
  \item A composition operation: given two arrows $f: A \rightarrow B$ and $g: B \rightarrow C$, there is an arrow $g \circ f: A \rightarrow C$.
  \item For every object $A$, an identity arrow $id_A: A \rightarrow A$.
\end{enumerate}

These data are required to satisfy the following axioms:

\begin{enumerate}
  \item \textbf{Associativity:} For arrows $f: A \rightarrow B$, $g: B \rightarrow C$, and $h: C \rightarrow D$, we have
  \[
  h \circ (g \circ f) = (h \circ g) \circ f.
  \]
  \item \textbf{Identity:} For every arrow $f: A \rightarrow B$, we have
  \[
  id_B \circ f = f \quad \text{and} \quad f \circ id_A = f.
  \]
\end{enumerate}

\begin{adjustwidth}{1cm}{1cm}
  \small
  \begin{quote}
For these reasons "set" turns out to have many meanings, so that the purported foundation of all of Mathematics upon set theory totters. But there are other possibilities. For example, the membership relation for sets can often be replaced by the composition operation for functions. This leads to an alternative foundation for Mathematics upon categoriesspecifically, on the category of all functions. Now much of Mathematics is dynamic, in that it deals with morphisms of an object L into another object of the same kind. Such morphisms (like functions) form categories, and so the approach via categories fits well with the objective of organizing and understanding Mathematics. That, in truth, should be the goal of a proper philosophy of Mathematics. \textcite[359]{maclane1986mathematics}
  \end{quote}
\end{adjustwidth}

\section{Other approaches}
\label{alternative-approaches}

As explained in the introduction, this text examines the mathematical significance of certain structures. In this section, I explore how this notion is addressed—often indirectly—within conceptual structuralism, a framework that has engaged extensively with CH. In the next section, I turn to how I would address the difficulties this conception raises for this issue.

From the perspective of conceptual structuralism, what can be evaluated is not the legitimacy of isolated facts and objects, but of the full structures or frameworks that are introduced in mathematical reasoning. For example, if the whole ZFC set theory universe is regarded as mathematically meaningful, then so is CH, since it is perfectly well-defined within that framework. In this regard, what Feferman most values, as mentioned in the introduction, is the degree to which such mathematical structures are conceptually clear and well-grounded. This emphasis is evident in the fifth of the ten theses he proposes to articulate conceptual structuralism:

\begin{adjustwidth}{1cm}{1cm}
    \small
    \begin{quote}
        Basic conceptions differ in their degree of clarity. One may speak of what is true in a given conception, but that notion of truth may only be partial. Truth in full is applicable only to completely clear conceptions. \parencite[p. 15]{Feferman2011a}
    \end{quote}
\end{adjustwidth}

Feferman argues that conceptions such as those required to formulate CH—namely, third-order arithmetic and stronger systems—suffer from an intrinsic lack of clarity. This is particularly evident in set theory, where, as he notes (\citeyear{Feferman2011a}, p.~1), “the concept of a set being arbitrary is vague or underdetermined, and there is no way to sharpen it without violating what it is supposed to be about”. For this reason, in his particular framework for understanding mathematics, he concludes that the CH cannot be assigned a definite truth value.

This argument has been subject to substantial criticism, most directly by Peter Koellner (\citeyear{Koellner2016} \& \citeyear{Koellner2022}). Koellner argues that Solomon Feferman’s notion of mathematical clarity is unstable and can be extended in such a way that even statements of arithmetic would be rendered indefinite—an outcome that Feferman himself explicitly rejects. I have examined this debate in \textcite{Blanco1}, where I argue that many of Koellner’s criticisms on this point are well taken. In that work, I also defended that Feferman’s position can be made more resilient by emphasizing an alternative interpretation of his conception—one that is closer to the notion of mathematical content that concerns me here. This interpretation focuses on what I call the \textit{applicability criterion}, namely the idea that mathematical entities exist through their systematic and consistent appearance in world pictures. 

Note: commitment to structures does ensure determinacy, while commitment to axiom systems does not, unless they are found to be a categorical set of axioms. Even if one could argue that the second order version of axioms in set theory is categorical, it is not true that they have been integrated to mathematical practice as the first order version has been. However, I wouldn't go as far as to state the same thing about third order arithmetic.

%---- Less on Feferman's applicability
Feferman adheres to this criterion with varying degrees of explicitness, though he emphasizes it quite generally. For example, it appears in Thesis 3 of conceptual structuralism (\citeyear{Feferman2011a}, p. 11):

\begin{adjustwidth}{1cm}{1cm}
\begin{quote}
\small
The basic conceptions of mathematics are of certain kinds of relatively simple ideal world-pictures which are not of objects in isolation but of structures, (...). They are communicated and understood prior to any axiomatics, indeed prior to any systematic logical development.
\end{quote}
\end{adjustwidth}

This passage emphasizes that the most basic and primitive mathematical structures are connected to the world we inhabit, in the sense that they help us understand and model it.\footnote{See section 2 in \textcite{Blanco1} for further discussion of this interpretation.} In another text, Feferman argues that the mathematics applicable to real-world physics can be derived from relatively simple systems such as Peano Arithmetic:

\begin{adjustwidth}{1cm}{1cm}
\small
\begin{quote}
I have come to conjecture that practically all scientifically applicable mathematics can be formalized in systems reducible to PA, or, as I have sloganized it in (1993): a little bit goes a long way. I have learned of a couple of cases (...) where it appears one must go beyond the resources of PA; but the physical theories that require such additional strength are rather speculative.
\end{quote} \hfill \parencite[p. 15]{Feferman1999}
\end{adjustwidth}

In conclusion, Feferman takes it to be important that mathematical structures are connected to real-world physics, at least insofar as they enable us to construct models and representations of it. However, he does not regard this connection as sufficient to establish the reality of certain mathematical conceptions—particularly those that figure in debates about the definiteness of the Continuum Hypothesis:

\begin{adjustwidth}{1cm}{1cm}
\small
\begin{quote}
Whether or not this thesis will hold for all future applications of mathematics to physical science, what this shows is that the argument from the success of physics to the reality of the set-theoretical conception of the real numbers is completely vitiated.
\end{quote} \hfill \parencite[p. 25]{Feferman2011a}
\end{adjustwidth}

However, a central difficulty with this emphasis is that it remains too restricted to direct applications in science, especially in physics, thereby overlooking—much like the indispensability argument—important aspects of mathematical practice. In particular, there are statements formulated within otherwise well-established frameworks, such as Fermat’s Last Theorem in PA, for which there is no known way of deciding them without appealing to stronger systems, such as finite-order arithmetic.\footnote{In this regard, see \textcite{McLarty2020}. I also discuss this problem in more detail in Section~\ref{section-CH-case}.} At the same time, it is crucial that mathematical entities be integrated into scientific practice in as simple and natural a way as possible. Some of the restrictions Feferman defends—such as limiting the treatment of the continuum to predicative, non-set-theoretical frameworks—do not seem to adequately account for this requirement.

Naturalists, such as \textcite{quine1953logical}, argued that mathematical entities and statements must in some sense be continuous with scientific theories. Thus, according to Quine, mathematical objects exist because of them being essential for the formulation of our best scientific theories, a view that came to be known as the \textit{indispensability argument}. However, \textcite{Maddy1997} defended that this argument ends up being incompatible with mathematical practice, and thus argued instead in favor of two fundamental complementary maxims: MAXIMIZE and UNIFY, maintaining that practitioners tend to make the mathematical universe as large and rich as possible, provided that consistency is preserved. 

\textbf{2. Mathematical Naturalism (W.V.O. Quine)}

Another strong parallel is Willard Van Orman Quine.

Mathematics is justified the same way as science: by its role in our best theories of the world.
No sharp boundary between math and empirical science.

Similarity:

Your $S_0$
 (structures grounded in physical applicability) is very Quinean.
Legitimacy comes from indispensability + usefulness.

Difference:

Quine emphasizes empirical indispensability, while your text adds internal criteria like unification and generativity.



\section{Conclusion}
\label{Conclusion}

It would be countersensical to exclude his subjective perspective from the investigation of the role that evidence plays in the selection of axioms. Admittedly, this contrasts with the concerns of ‘ordinary science’ where some standard of objectivity is presupposed and subjectivity is regarded as a disturbing element. \textcite[p. 104]{hauser2004}

Husserl's signitive vs intuitive distintion is crutial in mathematics. The comprehension of these two kind of consciousness acts is what he calls \textit{fulfillment}. Relation to sense vs reference.

We experience how in intuition the same objectuality is intuitively realized which was “merely meant” in the symbolic act, and that it becomes intuitable exactly as something determined in such a way as it was initially merely thought (merely meant). \textcite[Part 2, §8]{Husserl1913}

\begin{adjustwidth}{1cm}{1cm}
\small
\begin{quote}
Mathematics belongs to man, not to God. We are not interested in properties of the positive integers that have no descriptive meaning for finite man. When somebody proves a positive integer to exist, we shall show how to find it. If God has mathematics of its own that has to be done, let him do it himself. \textcite[p. 2]{Bishop1967}
\end{quote}
\end{adjustwidth}

\begin{adjustwidth}{1cm}{1cm}
  \small
  \begin{quote}
These considerations may lead us to say that $2^{\aleph_0} > \aleph_0$. That is to say: we can \emph{make} the considerations lead us to that. Or: we can say \emph{this} and give \emph{this} as our reason. But if we do say it---what are we to do next? In what practice is this proposition anchored? It is for the time being a piece of mathematical architecture which hangs in the air, and looks as if it were, let us say, an architrave, but not supported by anything and supporting nothing. \parencite[III, §35]{Wittgenstein1978}
  \end{quote}
\end{adjustwidth}

\begin{adjustwidth}{1cm}{1cm}
  \small
  \begin{quote}
Moreover, there are good reasons why Mathematicians do not usually present their proofs in fully formal style. It is because proofs are not only a means to certainty, but also a means to understanding. Behind each substantial formal proof there lies an idea, or perhaps several ideas. The idea, initially perhaps tenuous, explains why the result holds. The idea becomes Mathematics only when it can be formally expressed, but that expression must be so couched as to reveal the idea ; it will not do to bury the idea under the formalism. \textcite[378]{maclane1986mathematics}
  \end{quote}
\end{adjustwidth}

El árbol de las matemáticas es muy intrincado.

In my view, the foundationalist program in mathematics has failed to introduce itself in other areas of mathematical practice.

%=========================
\printbibliography


\end{document}
